% =====================================================================
%  Lecture 2026-11-05
%
%  Moved and reordered:
%   - "Lemma 1" (X x Y is Banach) is the proof of Proposition 20.4; it is
%     typeset there, in fragment 08, on instruction
%   - "Lemma 2" is placed before the closed graph theorem, whose proof
%     cites it
%
%  CORRECTED SOURCE ERRORS
%   [p1] resolvent "of X"  ->  of T at lambda
%   [p2] lemma on the domain of the resolvent: "T bounded or closed".
%        False for bounded T with a non-closed domain: for a dense proper
%        subspace D, T = 0 on D and lambda = -1, the inverse of
%        T - lambda I is the identity on D, bounded and densely defined,
%        so lambda is a regular value, but R_lambda(T) is defined on D
%        only  ->  "T in B(X) or T closed"; the source's proof covers the
%        closed case only
%
%  ADDITIONS
%   [p1] the domain Ran(T - lambda I) of R_lambda(T) named, and "densely
%        defined" glossed, since the definition of the spectrum uses them
%   [p1] integral-equation example: Example 8.10 cited for T in B(X)
%   [p2] shift example: Example 16.10 cited for S
%   [p1] Lemma 20.8 noted as the subspace case of Remark 13.3
%
%  REPORTED BUT UNMODIFIED GAPS
%   [p1] closed graph theorem: the implications (1) and (3) => (2) and
%        (2) and (3) => (1) are not proved, on instruction; the source
%        leaves them as exercises. The resolvent lemma uses the first.
%   [p1] closed graph theorem: G(T) a subspace, and P injective and
%        surjective, not stated
%   [p1] Lemma 20.8: "convergent implies Cauchy" and uniqueness of
%        limits used, not stated
%   [p2] the three parts of the spectrum are not said to partition it
%   [p2] shift example: injectivity of S, and the reason {y_1 = 0} is
%        not dense, not stated, on instruction
%   [p2] resolvent lemma: that T - lambda I and its inverse are closed
%        is asserted, on instruction; the case T in B(X) is not proved
%
%  OMITTED
%   none
% =====================================================================

% [p1]
\begin{lemma}\label{lem:closed-banach}
If $X$ is a Banach space and $Z \subset X$ is a subspace, then $Z$ is closed
if and only if $Z$ is a Banach space.
\end{lemma}

\begin{proof}
Assume $Z$ is closed, and let $(x_n) \subset Z$ be a Cauchy sequence. Then
$x_n$ has a limit $x$ in $X$, and $x$ must be in $Z$, since $Z$ is a closed
subspace.

Now assume $Z$ is a Banach space, and let $(x_n) \subset Z$ be a sequence
with $x_n \to x$. Then $(x_n)$ is Cauchy in $Z$, thus it converges in $Z$.
Therefore $x \in Z$, and $Z$ is closed.
\end{proof}

\begin{remark}
Lemma~\ref{lem:closed-banach} is the case of subspaces of the statement in
Remark~\ref{rem:minimizing-closed}.
\end{remark}

\begin{theorem}[Closed graph theorem]\label{thm:closed-graph}
Let $X$ and $Y$ be Banach spaces, and let
$T \colon \operatorname{Dom} T \subset X \to Y$ be linear, with
$\operatorname{Dom} T$ a subspace. Then any two of the following properties
imply the third:
\begin{enumerate}
  \item $T$ is closed;
  \item $\operatorname{Dom} T$ is closed;
  \item $T$ is bounded, $T \in \Bop{\operatorname{Dom} T}{Y}$.
\end{enumerate}
\end{theorem}

\begin{proof}
The implications (1) and (3) $\Rightarrow$ (2), and (2) and (3)
$\Rightarrow$ (1), are not proved here. The nontrivial case is (1) and (2)
$\Rightarrow$ (3).

If $X$ and $Y$ are Banach spaces, then so is $X \times Y$ with the norm
$\norm{(x,y)} = \norm{x} + \norm{y}$, by
Proposition~\ref{prop:product-space}. So $G(T) \subset X \times Y$ is a
closed subspace of a Banach space, and is therefore a Banach space, by
Lemma~\ref{lem:closed-banach}. Now consider the projection
\[
  P \colon G(T) \to \operatorname{Dom} T, \qquad (x, Tx) \mapsto x .
\]
$P$ is clearly linear, and it is bounded, since
$\norm{x} \le \norm{(x, Tx)} = \norm{x} + \norm{Tx}$, with constant $1$, so
$\norm{P} \le 1$. $P$ is also a bijection, thus $P^{-1}$ exists:
$P^{-1} \colon x \mapsto (x, Tx)$. $G(T)$ is a Banach space, and
$\operatorname{Dom} T$ is closed in a Banach space, therefore also a Banach
space. So by the bounded inverse theorem,
Corollary~\ref{cor:bounded-inverse}, $P^{-1}$ is bounded. So there is a $c$
with
\[
  \norm{P^{-1}(x)} \le c\norm{x}
  \iff \norm{x} + \norm{Tx} \le c\norm{x}
  \iff \norm{Tx} \le (c - 1)\norm{x} ,
\]
and $T$ is bounded.
\end{proof}

\begin{remark}
The class of closed linear operators is introduced because in applications
we need differential operators, which are not bounded, but we need some
topological substitute for boundedness.
\end{remark}

% [p1]
\section{Spectral theory}

\begin{remark}
The spaces considered in spectral theory are always complex.
\end{remark}

\begin{definition}[Resolvent]\label{def:resolvent}
Let $X$ be a Banach space, $T \colon \operatorname{Dom} T \subset X \to X$
linear with $\operatorname{Dom} T$ a subspace, and $\lambda \in \C$. Then
$R_\lambda(T) = (T - \lambda I)^{-1}$, if it exists, that is, if
$T - \lambda I$ is injective, is called the \emph{resolvent} of $T$ at
$\lambda$. It is defined on $\operatorname{Ran}(T - \lambda I)$, the range of
$T - \lambda I$.
\end{definition}

\begin{example}
Let $X = \Cont{a,b}$ and let $y \in X$ be fixed. We want to find $x \in X$
such that
\[
  \int_a^b k(s,t)\, x(t) \dd t - \lambda x(s) = y(s) ,
\]
where the kernel $k$ is in $C([a,b] \times [a,b])$. We can write
$T \colon X \to X$, $(Tx)(s) = \int_a^b k(s,t)\, x(t) \dd t$; then
$T \in \Bopone{X}$, by Example~\ref{ex:integral-kernel} with $[a,b]$ in place
of $[0,1]$. We want to solve $Tx - \lambda x = y$, and hopefully
$x = R_\lambda(T) y$.
\end{example}

\begin{definition}[Spectrum]\label{def:spectrum}
Let $X$ be a Banach space and $T \colon \operatorname{Dom} T \subset X \to X$
linear. We say that $\lambda \in \C$ is a \emph{regular value} of $T$ if
$R_\lambda(T) = (T - \lambda I)^{-1}$ exists, is bounded and is densely
defined, that is, its domain $\operatorname{Ran}(T - \lambda I)$ is dense in
$X$. If $\lambda$ is not a regular value, then we call it a value in the
\emph{spectrum} of $T$.
\[
  \rho(T) := \set{\lambda \in \C \mid \lambda \text{ is a regular value of }
  T}
\]
is the \emph{resolvent set}, and $\sigma(T) = \C \setminus \rho(T)$ is the
\emph{spectrum} of $T$.
\end{definition}

% [p2]
\begin{definition}[Parts of the spectrum]\label{def:spectrum-parts}
A value $\lambda \in \sigma(T)$ is in
\begin{enumerate}
  \item the \emph{point spectrum}, if $(T - \lambda I)^{-1}$ does not exist,
    that is, $T - \lambda I$ is not injective;
  \item the \emph{continuous spectrum}, if $(T - \lambda I)^{-1}$ exists and
    is densely defined, but not bounded;
  \item the \emph{residual spectrum}, if $(T - \lambda I)^{-1}$ exists, but
    is not densely defined.
\end{enumerate}
\end{definition}

\begin{example}
Let $S$ be the shift operator of Example~\ref{ex:shift},
$S \colon \lp{2} \to \lp{2}$,
$(x_1, x_2, x_3, \dots) \mapsto (0, x_1, x_2, x_3, \dots)$. We claim that
$0 \in \sigma(S)$ is in the residual spectrum. For $\lambda = 0$,
$(S - \lambda I)^{-1}$ exists; it is
\[
  S^{-1} \colon (y_1, y_2, y_3, \dots) \longmapsto (y_2, y_3, \dots) ,
  \qquad
  \operatorname{Dom}(S^{-1}) = \operatorname{Ran} S
  = \set{y \in \lp{2} \mid y_1 = 0} ,
\]
with $S^{-1} \colon \operatorname{Dom}(S^{-1}) \to \lp{2}$. So $S^{-1}$ is not
densely defined, and $0$ is in the residual spectrum.
\end{example}

\begin{lemma}\label{lem:resolvent-domain}
Let $X$ be a Banach space, $T \colon \operatorname{Dom} T \subset X \to X$
linear, and either $T \in \Bopone{X}$ or $T$ closed, and let
$\lambda \in \rho(T)$. Then $R_\lambda(T)$ is defined on the whole space
$X$.
\end{lemma}

\begin{proof}
If $T$ is closed, then $T - \lambda I$ is closed, and
$(T - \lambda I)^{-1}$ is also closed. But $(T - \lambda I)^{-1}$ is bounded
and densely defined: $\operatorname{Dom}((T - \lambda I)^{-1})
= \operatorname{Ran}(T - \lambda I)$ is dense. This implies that
$\operatorname{Dom}((T - \lambda I)^{-1})$ is closed, by the closed graph
theorem, Theorem~\ref{thm:closed-graph}. So
$\operatorname{Dom}((T - \lambda I)^{-1}) = X$.
\end{proof}
